\documentclass[11pt,openany]{article}

\input{algebra-note-preamble}
\input{tcolorbox}
\input{theorem}
\input{algebra-note-commands}
\usepackage{tikz}
\usepackage{tikz-3dplot}
\usepackage{amsmath}
\usepackage{tikz-cd}

\setstretch{1.25}
\begin{document}
\pagenumbering{arabic}
\begin{center}
	\huge\textbf{Quotient Rings, Closed Sets,\\ and Maximal Spectra}\\
	\vspace{0.5em}
	\large{Ji, Yong-hyeon}\\
	\vspace{0.5em}
	\normalsize{\today}\\
\end{center}

Let:
\begin{itemize}
	\item \( A = C[x] \), the ring of polynomials over \( \mathbb{C} \) (the field of complex numbers),
	\item \( S = \{1, x, x^2, \dots\} \), the multiplicative set generated by \( x \), consisting of all powers of \( x \), i.e., \( S = \{x^n \mid n \geq 0\} \),
	\item Then, \( S^{-1}A = C[x, x^{-1}] \), the localization of \( C[x] \) at \( S \), which gives the ring of Laurent polynomials (polynomials in \( x \) and \( x^{-1} \) with coefficients in \( \mathbb{C} \)).
\end{itemize}

Define a ring map:
\[
f: C[x] \to \mathbb{C}, \quad f(p(x)) = p(1),
\]
i.e., \( f \) evaluates the polynomial at \( x = 1 \). Note that \( f(x) = 1 \), which is a unit in \( \mathbb{C} \).

\subsection*{Claim}
There exists a unique homomorphism \( g : S^{-1}A \to \mathbb{C} \) such that the following diagram commutes:
\[
\begin{aligned}
	& C[x] \xrightarrow{f} \mathbb{C} \\
	& \downarrow \\
	& S^{-1}C[x] \xrightarrow{g} \mathbb{C}
\end{aligned}
\]

\subsection*{Verification}

\paragraph{Existence of \( g \):} In \( S^{-1}C[x] \), every element is of the form \( \frac{p(x)}{x^n} \), where \( p(x) \in C[x] \) and \( n \geq 0 \). Define:
\[
g\left(\frac{p(x)}{x^n}\right) = \frac{f(p(x))}{f(x)^n}.
\]
Since \( f(x) = 1 \), this simplifies to:
\[
g\left(\frac{p(x)}{x^n}\right) = \frac{p(1)}{1^n} = p(1).
\]

\paragraph{Commutativity of the Diagram:} For any \( a \in C[x] \), its image under the map \( A \to S^{-1}A \to \mathbb{C} \) equals its direct image under \( f \). For example, let \( p(x) = x^2 + 3x + 2 \):
\[
f(p(x)) = p(1) = 1^2 + 3(1) + 2 = 6,
\]
and
\[
g\left(\frac{p(x)}{1}\right) = p(1) = 6.
\]
Thus, the diagram commutes.

\paragraph{Uniqueness of \( g \):} Any other homomorphism \( g' \) must agree with \( f \) on \( C[x] \) and satisfy \( g'(x) = f(x) = 1 \). This uniquely determines \( g \).

\subsection*{Conclusion}
This example illustrates the proposition in the context of \( C[x] \). The ring \( S^{-1}C[x] \) (the Laurent polynomials \( C[x, x^{-1}] \)) allows us to extend the map \( f : C[x] \to \mathbb{C} \) uniquely to \( g : S^{-1}C[x] \to \mathbb{C} \). The localization \( S^{-1}C[x] \) captures the multiplicative inverses of elements in \( S \), ensuring that \( g \) is well-defined and commutative.

%\section{Visualization of the Example}

%The following diagram illustrates the relationship between ( C[x] ), ( S^{-1}C[x] ), and ( \mathbb{C} ):


%\begin{tikzpicture}[node distance=3cm, auto, thick]
%	
%	% Nodes
%	\node (A) at (0, 2) {( C[x] )};
%	\node (B) at (3, 2) {( \mathbb{C} )};
%	\node (C) at (0, 0) {( S^{-1}C[x] )};
%	\node (D) at (3, 0) {( \mathbb{C} )};
%	
%	% Arrows
%	\draw[->] (A) -- node[above] {( f(p(x)) = p(1) )} (B);
%	\draw[->] (A) -- node[left] {( \text{localization} )} (C);
%	\draw[->] (C) -- node[below] {( g(\frac{p(x)}{x^n}) = p(1) )} (D);
%	\draw[->] (B) -- node[right] {( \text{id} )} (D);
%	
%	% Labels for sets
%	\node[align=center] at (-1.5, 2) {Ring \ ( C[x] )};
%	\node[align=center] at (-1.5, 0) {Localized \ ( S^{-1}C[x] )};
%	\node[align=center] at (4.5, 2) {( \mathbb{C} )};
%	\node[align=center] at (4.5, 0) {( \mathbb{C} )};
%	
%\end{tikzpicture}

%\subsection{Explanation}
%\begin{itemize}
%	\item ( C[x] ) is the ring of polynomials over ( \mathbb{C} ).
%	\item ( S^{-1}C[x] ) is the localization of ( C[x] ) at the multiplicative set ( S = {x^n \mid n \geq 0} ), giving the ring of Laurent polynomials.
%	\item ( f: C[x] \to \mathbb{C} ) is a ring homomorphism defined by ( f(p(x)) = p(1) ).
%	\item ( g: S^{-1}C[x] \to \mathbb{C} ) extends ( f ) to the localized ring.
%	\item The diagram commutes, meaning the same result is achieved whether we apply ( f ) directly or first localize via ( g ).
%\end{itemize}

\newpage
\begin{tikzpicture}[>=stealth]
	\matrix[matrix of math nodes,row sep=2.5em,column sep=4em,ampersand replacement=\&]
	{
		C[x] \& \mathbb{C} \\
		S^{-1}C[x] = C[x]_{(x-a)} \& \mathbb{C} \\
	};
	\path[->]
	(1-1) edge node[above] {\( f(p(x)) = p(a) \)} (1-2)
	(1-1) edge node[left] {\( \text{localization} \)} (2-1)
	(2-1) edge node[below] {\( g \)} (2-2)
	(1-2) edge[->] node[right] {\( \text{id} \)} (2-2);
\end{tikzpicture}

Here:
- \( f(p(x)) = p(a) \): Evaluates the polynomial at \( x = a \).
- \( g\left(\frac{p(x)}{q(x)}\right) = \frac{p(a)}{q(a)} \) in \( S^{-1}C[x] = C[x]_{(x-a)} \), focusing on the behavior near \( x = a \).

\subsection*{2. Localization at Complement of \( x-a \)}

\begin{tikzpicture}[>=stealth]
	\matrix[matrix of math nodes,row sep=2.5em,column sep=4em,ampersand replacement=\&]
	{
		C[x] \& \mathbb{C} \\
		S^{-1}C[x] = C[x]_{(x-a)} \& \mathbb{C} \\
	};
	\path[->]
	(1-1) edge node[above] {\( f(p(x)) = p(a) \)} (1-2)
	(1-1) edge node[left] {\( \text{localization} \)} (2-1)
	(2-1) edge node[below] {\( g \)} (2-2)
	(1-2) edge[->] node[right] {\( \text{id} \)} (2-2);
\end{tikzpicture}

Here:
- \( S = C[x] \setminus (x-a) \), where \( S \) includes all polynomials not divisible by \( x-a \).
- \( f(p(x)) = p(a) \) remains the evaluation at \( x = a \).
- \( g \) extends \( f \) to the localization by allowing division by elements in \( S \).


\subsection*{3. Localization at Complement of Zero (\( S = C[x] \setminus \{0\} \))}

\begin{tikzpicture}[>=stealth]
%	\matrix[matrix of math nodes,row sep=2.5em,column sep=4em,ampersand replacement=\&]
%	{
%		$C[x]$ \& $C(x)$ \\
%		$S^{-1}C[x] = C(x$) \&$ C(x)$ \\
%	};
%	\path[->]
%	(1-1) edge node[above] {\( f(p(x)) \)} (1-2)
%	(1-1) edge node[left] {\( \text{localization} \)} (2-1)
%	(2-1) edge node[below] {\( g = \text{id} \)} (2-2)
%	(1-2) edge[->] node[right] {\( \text{id} \)} (2-2);
\end{tikzpicture}
\newpage
\begin{proposition}
	Let $f:A\to B$ be a ring homomorphism and $S\subseteq A$ be a multiplicative set such that $f(s)$ is a unit in $B$ for all $s\in S$. Then there is a unique homomorphism $g:S^{-1}A\to B$ such that the following diagram commutes:
%	\begin{figure}[h!]\centering
%	\begin{center}
%		\begin{tikzcd}
%			A \arrow[rrddd] \arrow[rrrr, "f"] &  &                             &  & B \\
%			&  &                             &  &   \\
%			&  &                             &  &   \\
%			&  & S^{-1}A \arrow[rruuu, "g"'] &  &  
%		\end{tikzcd}
%	\end{center}
%	\end{figure}% https://tikzcd.yichuanshen.de/#N4Igdg9gJgpgziAXAbVABwnAlgFyxMJZABgBpiBdUkANwEMAbAVxiRAEEQBfU9TXfIRQAWclVqMWbAELdeIDNjwEiAJlIBmcfWatEIAMoA9YAFoAjF05dxMKAHN4RUADMAThAC2SMiBwQkVR5XD29EX38kc2odKX0XORCvQOpIxGiJXTZ7EGoGOgAjGAYABX5lIRA3LHsACxxuCi4gA
\end{proposition}
\begin{center}
	\begin{tikzcd}
		A \arrow[rrddd] \arrow[rrrr, "f"] &  &                             &  & B \\
		&  &                             &  &   \\
		&  &                             &  &   \\
		&  & S^{-1}A \arrow[rruuu, "g"'] &  &  
	\end{tikzcd}
\end{center}

\end{document}
